% ថ្នាក់ទី ១១ (វិទ្យាសាស្ត្រ) · ជំពូកទី ៣ · មេរៀនទី ៥ · សៀវភៅទំព័រ ១៤២–១៦០
\chapter{សមីការនិងវិសមីការត្រីកោណមាត្រផ្សេងទៀត}

\begin{objBox}[វត្ថុបំណង]
\begin{enumerate}[1]
  \item ដោះស្រាយសមីការត្រីកោណមាត្រ
  \item ដោះស្រាយវិសមីការត្រីកោណមាត្រ ។
\end{enumerate}
\end{objBox}

យើងបានសិក្សារួចមកហើយពីសមីការនិងវិសមីការត្រីកោណមាត្រងាយៗ នៅក្នុងសៀវភៅគណិតវិទ្យាកម្រិតមូលដ្ឋាន ។ ក្នុងផ្នែកនេះយើងនឹងសិក្សាពីសមីការនិងវិសមីការត្រីកោណមាត្រដែលមានទម្រង់ផ្សេងៗទៀត ។

\section{សមីការត្រីកោណមាត្រ}

\subsection{សមីការដឺក្រេទី 1 ធៀបនឹង \texorpdfstring{$\sin x$}{sin x} និង \texorpdfstring{$\cos x$}{cos x}}

ជាសមីការត្រីកោណមាត្រដែលក្រោយពីបង្រួមហើយមានរាង $a\cos x + b\sin x = c$ ។ ដំណោះស្រាយសមីការនេះមានពីររបៀប ៖

\textbf{របៀបទី ១} ៖
\begin{itemize}
  \item បំប្លែងសមីការនេះជារាង $\cos(x - \theta) = \dfrac{c}{r}$ ដែល $r = \sqrt{a^2 + b^2}$, $(r \ge 0)$ និង $\cos\theta = \dfrac{a}{r}$, $\sin\theta = \dfrac{b}{r}$ ។ ដោយ $\cos\theta = \dfrac{a}{r}$ នាំឱ្យ $a = r\cos\theta$, $\sin\theta = \dfrac{b}{r}$ នាំឱ្យ $b = r\sin\theta$ ។

  សមីការ $a\cos x + b\sin x = r\cos\theta\cos x + r\sin\theta\sin x = c$
  \[
\begin{gathered}
r(\cos x\cos\theta + \sin x\sin\theta) = c \quad \text{សមមូល}\\
r\cos(x - \theta) = c \quad \text{ឬ}\\
\cos(x - \theta) = \frac{c}{r} \text{ ។}
\end{gathered}
\]
  \item ដោះស្រាយសមីការ $\cos(x - \theta) = \dfrac{c}{r}$ តាមរូបមន្ត $\cos x = \cos\alpha$ ។
\end{itemize}

\begin{example}[ឧទាហរណ៍ ១]
ដោះស្រាយសមីការ $\sin x + \sqrt{3}\cos x = 0$ \ (1) ។

គេបាន $r = \sqrt{1 + 3} = 2$, $\cos\theta = \dfrac{a}{r} = \dfrac{1}{2}$, $\sin\theta = \dfrac{b}{r} = \dfrac{\sqrt{3}}{2}$ នាំឱ្យ $\theta = \dfrac{\pi}{3}$ ។
\[
\begin{aligned}
\sin x + \sqrt{3}\cos x &= 2\cos\frac{\pi}{3}\sin x + 2\sin\frac{\pi}{3}\cos x = 0\\
2\left(\sin x\cos\frac{\pi}{3} + \cos x\sin\frac{\pi}{3}\right) &= 0\\
2\sin\left(x + \frac{\pi}{3}\right) = 0 \quad &\text{ឬ} \quad \sin\left(x + \frac{\pi}{3}\right) = 0
\end{aligned}
\]
សមីការ (1) មានចម្លើយ $\sin\left(x + \dfrac{\pi}{3}\right) = 0$ នាំឱ្យ $x + \dfrac{\pi}{3} = k\pi$ ។

ដូច្នេះ $x = -\dfrac{\pi}{3} + k\pi$, $(k \in \mathbb{Z})$ ។ ចំនួនចុងធ្នូចម្លើយមានពីរ ។
\end{example}

\textbf{របៀបទី ២} ៖ គណនា $\sin x$ និង $\cos x$ ជាអនុគមន៍នៃ $\tan\dfrac{x}{2}$ ដោយ $x \ne \pi + 2k\pi$, $(k \in \mathbb{Z})$ ។

គេជំនួស $\cos x = \dfrac{1 - t^2}{1 + t^2}$ និង $\sin x = \dfrac{2t}{1 + t^2}$ នោះសមីការ $a\cos x + b\sin x = c$ ទៅជាសមីការដឺក្រេទី $2$ ដែលមាន $t$ ជាអញ្ញាត ៖ $(c + a)t^2 - 2bt + c - a = 0$ ។

ដោះស្រាយសមីការដឺក្រេទី $2$ រកឫសរួចដោះស្រាយសមីការតាមរូបមន្ត $\tan\dfrac{x}{2} = \tan\alpha$ ។

\begin{note}[សម្គាល់]
វិធីនេះអាចបាត់ឫសត្រង់ $x = \pi + 2k\pi$ ។ ដូច្នេះ ដំបូងគេត្រូវពិនិត្យមើលថាតើ $x = \pi + 2k\pi$ ជាឫសរបស់សមីការ $a\cos x + b\sin x = c$ ឬទេ?
\end{note}

\begin{example}[ឧទាហរណ៍ ២]
ដោះស្រាយសមីការ $\sin x + 7\cos x = 5$ ។

តាង $t = \tan\dfrac{x}{2}$ $(x \ne \pi + 2k\pi)$ នោះសមីការសរសេរ $\dfrac{2t}{1 + t^2} + 7\left(\dfrac{1 - t^2}{1 + t^2}\right) = 5$
\[ 2t + 7 - 7t^2 = 5 + 5t^2 \quad \text{ឬ} \quad 6t^2 - t - 1 = 0 \]
សមីការមានឫស $t_1 = \dfrac{1}{2}$, $t_2 = -\dfrac{1}{3}$ ។
\begin{itemize}
  \item បើ $t_1 = \dfrac{1}{2}$ គេបាន $\tan\dfrac{x}{2} = \dfrac{1}{2}$ នោះមាន $\alpha_1 \in \mathbb{R}$ ដែល $\tan\alpha_1 = \dfrac{1}{2}$ ។ គេបាន $\tan\dfrac{x}{2} = \tan\alpha_1$ នាំឱ្យ $\dfrac{x}{2} = \alpha_1 + k\pi$ ឬ $x = 2\alpha_1 + 2k\pi$, $(k \in \mathbb{Z})$ ។
  \item បើ $t_2 = -\dfrac{1}{3}$ គេបាន $\tan\dfrac{x}{2} = -\dfrac{1}{3}$ នោះមាន $\alpha_2 \in \mathbb{R}$ ដែល $\tan\alpha_2 = -\dfrac{1}{3}$ ។ គេបាន $\tan\dfrac{x}{2} = \tan\alpha_2$ នាំឱ្យ $\dfrac{x}{2} = \alpha_2 + k\pi$ ឬ $x = 2\alpha_2 + 2k\pi$, $(k \in \mathbb{Z})$ ។
\end{itemize}
ដូច្នេះ សមីការមានចម្លើយ ៖ $x = 2\alpha_1 + 2k\pi$, $x = 2\alpha_2 + 2k\pi$, $(k \in \mathbb{Z})$ ។
\end{example}

\begin{example}[លំហាត់គំរូ ១]
ដោះស្រាយសមីការ $\sqrt{3}\sin 2x + \cos 2x = \sqrt{2}$ \ (1) ។
\end{example}
\ans
គេបាន $\cos 2x + \sqrt{3}\sin 2x = \sqrt{2}$ ហើយ $r = \sqrt{3 + 1} = 2$, $\cos\theta = \dfrac{1}{2}$, $\sin\theta = \dfrac{b}{r} = \dfrac{\sqrt{3}}{2}$ ។

សមីការសរសេរ $2\left(\dfrac{1}{2}\cos 2x + \dfrac{\sqrt{3}}{2}\sin 2x\right) = \sqrt{2}$ ឬ $\cos\left(2x - \dfrac{\pi}{3}\right) = \dfrac{\sqrt{2}}{2} = \cos\dfrac{\pi}{4}$ ។

សមីការ (1) មានចម្លើយ ៖
\[
\begin{gathered}
2x - \frac{\pi}{3} = \frac{\pi}{4} + 2k\pi, \ x = \frac{7\pi}{24} + k\pi \quad \text{និង}\\
2x - \frac{\pi}{3} = -\frac{\pi}{4} + 2k\pi, \ x = \frac{\pi}{24} + k\pi, \ (k \in \mathbb{Z})
\end{gathered}
\]
ចំនួនចុងធ្នូចម្លើយមាន $4$ ។

\begin{pro}
ដោះស្រាយសមីការ
\begin{tasks}(2)
  \task $\sin x + \sqrt{3}\cos x = \sqrt{2}$
  \task $\cos 2x - \sin 2x = -1$
  \task $\cos x + \sqrt{3}\sin x = 1$
  \task $\cos x - \sqrt{3}\sin x = -\sqrt{2}$ ។
\end{tasks}
\ans
\begin{enumerate}[a]
  \item $x = \dfrac{5\pi}{12} + 2k\pi$, $x = -\dfrac{\pi}{12} + 2k\pi$, $(k \in \mathbb{Z})$
  \item $x = \dfrac{\pi}{2} + k\pi$, $x = \dfrac{\pi}{4} + k\pi$, $(k \in \mathbb{Z})$
  \item $x = \dfrac{2\pi}{3} + 2k\pi$, $x = 2k\pi$, $(k \in \mathbb{Z})$
  \item $x = \dfrac{5\pi}{12} + 2k\pi$, $x = -\dfrac{13\pi}{12} + 2k\pi$, $(k \in \mathbb{Z})$ ។
\end{enumerate}
\end{pro}

\subsection{សមីការដឺក្រេទី 2 ធៀបនឹងអនុគមន៍ត្រីកោណមាត្រនៃចំនួនពិត \texorpdfstring{$x$}{x}}

សមីការទាំងនោះមានរាង ៖ $a\cos^2 x + b\cos x + c = 0$, $a\sin^2 x + b\sin x + c = 0$, $a\tan^2 x + b\tan x + c = 0$ និង $a\cot^2 x + b\cot x + c = 0$ ។ ដើម្បីដោះស្រាយសមីការបែបនេះ គេត្រូវ ៖
\begin{enumerate}[1]
  \item បំប្លែងសមីការឱ្យទៅជាសមីការដឺក្រេទី $2$ ដូចខាងលើសិន ។
  \item តាងអនុគមន៍ត្រីកោណមាត្រនៃចំនួនពិត $x$ ដោយ $X$ ឬ $Y$ ឬ $t \ldots$ ។
  \item ដាក់លក្ខខណ្ឌទៅតាមអនុគមន៍ត្រីកោណមាត្រដែលឱ្យ ។
  \item ដោះស្រាយសមីការដឺក្រេទី $2$ ដែលមានអញ្ញាត $X$ ឬ $Y$ ឬ $t \ldots$ ។
  \item យកឫសដែលបានមកជំនួសជាអនុគមន៍ត្រីកោណមាត្រវិញ រួចផ្ទៀងផ្ទាត់លក្ខខណ្ឌខាងលើ ហើយដោះស្រាយតាមរូបមន្តសមីការត្រីកោណមាត្រ ។
\end{enumerate}

\begin{example}[ឧទាហរណ៍ ១]
ដោះស្រាយសមីការ $4 - \cos 2x - 7\sin x = 0$ ។

ដោយ $\cos 2x = 1 - 2\sin^2 x$ សមីការអាចសរសេរ $2\sin^2 x - 7\sin x + 3 = 0$ \ (1) ។ តាង $t = \sin x$ ដែលមានលក្ខខណ្ឌ $-1 \le t \le 1$ ។

សមីការ (1) អាចសរសេរ $2t^2 - 7t + 3 = 0$, $\Delta = 25$, $t_1 = \dfrac{1}{2}$, $t_2 = 3$ មិនផ្ទៀងផ្ទាត់លក្ខខណ្ឌខាងលើ ។

ដូច្នេះ $t_1 = \sin x = \dfrac{1}{2} = \sin\dfrac{\pi}{6}$ យើងបាន $x_1 = \dfrac{\pi}{6} + 2k\pi$, $x_2 = \dfrac{5\pi}{6} + 2k\pi$, $(k \in \mathbb{Z})$ ។
\end{example}

\begin{example}[ឧទាហរណ៍ ២]
ដោះស្រាយសមីការ $\tan^2\dfrac{x}{2} - (1 - \sqrt{3})\tan\dfrac{x}{2} - \sqrt{3} = 0$ ។

គេតាង $\tan\dfrac{x}{2} = t$ សមីការ (1) សរសេរ $t^2 - (1 - \sqrt{3})t - \sqrt{3} = 0$ ។ ដោះស្រាយសមីការនេះ គេបាន $t_1 = -1$, $t_2 = \sqrt{3}$
\[
\begin{gathered}
t_1 = \tan\frac{x}{2} = -1 = \tan\left(-\frac{\pi}{4}\right),\\
t_2 = \tan\frac{x}{2} = \sqrt{3} = \tan\frac{\pi}{3} \text{ ។}
\end{gathered}
\]
ដូច្នេះ $\tan\dfrac{x}{2} = \tan\left(-\dfrac{\pi}{4}\right)$, $\dfrac{x}{2} = -\dfrac{\pi}{4} + k\pi$, $x = -\dfrac{\pi}{2} + 2k\pi$ ។

$\tan\dfrac{x}{2} = \tan\dfrac{\pi}{3}$, $\dfrac{x}{2} = \dfrac{\pi}{3} + k\pi$, $x = \dfrac{2\pi}{3} + 2k\pi$, $(k \in \mathbb{Z})$ ។

ដូច្នេះ សមីការមានចម្លើយ ៖ $x = -\dfrac{\pi}{2} + 2k\pi$ និង $x = \dfrac{2\pi}{3} + 2k\pi$, $(k \in \mathbb{Z})$ ។
\end{example}

\begin{example}[លំហាត់គំរូ ២]
ដោះស្រាយសមីការ $\dfrac{\cos x(2\sin x + 3\sqrt{2}) - 2\cos^2 x - 1}{1 + \sin 2x} = 1$ ។
\end{example}
\ans
លក្ខខណ្ឌ $1 + \sin 2x \ne 0$, $\sin 2x \ne -1$, $2x \ne -\dfrac{\pi}{2} + 2k\pi$, $x \ne -\dfrac{\pi}{4} + k\pi$, $(k \in \mathbb{Z})$ ។

សមីការសរសេរ $\cos x(2\sin x + 3\sqrt{2}) - 2\cos^2 x - 1 = 1 + \sin 2x$
\[
\begin{gathered}
2\sin x\cos x + 3\sqrt{2}\cos x - 2\cos^2 x - 1 = 1 + \sin 2x,\\
2\cos^2 x - 3\sqrt{2}\cos x + 2 = 0
\end{gathered}
\]
តាង $t = \cos x$ សមីការទៅជា $2t^2 - 3\sqrt{2}t + 2 = 0$ ។ ដោះស្រាយសមីការនេះ គេបាន $t_1 = \sqrt{2}$, $t_2 = \dfrac{1}{\sqrt{2}}$ ។

$t_1 = \cos x = \sqrt{2}$ (មិនយក) ។ $t_2 = \cos x = \dfrac{1}{\sqrt{2}} = \cos\dfrac{\pi}{4}$ ។ ដូច្នេះ $\cos x = \cos\dfrac{\pi}{4}$ នាំឱ្យ $x = \pm\dfrac{\pi}{4} + 2k\pi$, $(k \in \mathbb{Z})$ ។

តាមលក្ខខណ្ឌខាងលើសមីការមានឫស $x = \dfrac{\pi}{4} + 2k\pi$, $(k \in \mathbb{Z})$ ។

\begin{pro}
ដោះស្រាយសមីការ
\begin{tasks}(2)
  \task $2\cos^2 x - 3\sqrt{2}\cos x + 2 = 0$
  \task $\dfrac{1}{\sin^2 x} = \cot x + 3$
  \task $5\tan^2 y + 5\tan y = 2(1 + \tan^2 y)$
  \task $8\sin^2 x - 6\sin x = 5$ ។
\end{tasks}
\ans
\begin{enumerate}[a]
  \item $x = \pm\dfrac{\pi}{4} + 2k\pi$, $(k \in \mathbb{Z})$
  \item $\cot x = -1$ ឬ $\cot x = 2$ ៖ $x = \dfrac{3\pi}{4} + k\pi$ ឬ $x = \alpha + k\pi$ ដែល $\cot\alpha = 2$, $(k \in \mathbb{Z})$
  \item $\tan y = \dfrac{1}{3}$ ឬ $\tan y = -2$ ៖ $y = \alpha + k\pi$ ឬ $y = \beta + k\pi$ ដែល $\tan\alpha = \dfrac{1}{3}$, $\tan\beta = -2$
  \item $\sin x = -\dfrac{1}{2}$ ៖ $x = -\dfrac{\pi}{6} + 2k\pi$ ឬ $x = \dfrac{7\pi}{6} + 2k\pi$, $(k \in \mathbb{Z})$ ។
\end{enumerate}
\end{pro}

\section{សមីការដែលមានរាង \texorpdfstring{$a\sin^2 x + b\sin x\cos x + c\cos^2 x = d$}{a sin²x + b sin x cos x + c cos²x = d}}

ដើម្បីដោះស្រាយសមីការនេះ គេត្រូវចែកអង្គទាំងពីរនៃសមីការនឹង $\cos^2 x \ne 0$, $x \ne \dfrac{\pi}{2} + k\pi$ ។ (បើ $d \ne 0$ គេសរសេរ $d = d(\sin^2 x + \cos^2 x)$ ជាមុន ។) គេបានសមីការ $a\tan^2 x + b\tan x + c = 0$ រួចដោះស្រាយតាមសមីការដឺក្រេទី $2$ ដូចខាងលើ ។

\begin{example}[ឧទាហរណ៍ ១]
ដោះស្រាយសមីការ $3\sin^2 x + 3\sin x\cos x - 6\cos^2 x = 0$ ។

ចែកអង្គទាំងពីរនៃសមីការនឹង $\cos^2 x$ សមីការសរសេរ $3\tan^2 x + 3\tan x - 6 = 0$ ។ តាង $t = \tan x$ គេបាន $3t^2 + 3t - 6 = 0$ សមីការនេះមានឫសពីរគឺ $t_1 = 1$, $t_2 = -2$ ។

$t_1 = \tan x = 1$ គេបាន $\tan x = \tan\dfrac{\pi}{4}$, $x = \dfrac{\pi}{4} + k\pi$, $(k \in \mathbb{Z})$ ។

$t_2 = -2$ តាង $\tan\alpha = -2$, $x = \alpha + k\pi$, $(k \in \mathbb{Z})$ ។ សមីការមានចម្លើយ $x = \dfrac{\pi}{4} + k\pi$ និង $x = \alpha + k\pi$, $(k \in \mathbb{Z})$ ។
\end{example}

\begin{example}[ឧទាហរណ៍ ២]
ដោះស្រាយសមីការ $\dfrac{\sin^4\frac{x}{2} + \cos^4\frac{x}{2}}{1 - \sin x} - \tan^2 x\sin x = \dfrac{1 + \sin x}{2} + \tan^2 x$ ។

លក្ខខណ្ឌ $\sin x \ne 1$, $\cos x \ne 0$, $x \ne \dfrac{\pi}{2} + k\pi$ ។ សមីការអាចសរសេរ
\[
\begin{aligned}
&\frac{\left(\sin^2\frac{x}{2} + \cos^2\frac{x}{2}\right)^2 - 2\sin^2\frac{x}{2}\cos^2\frac{x}{2}}{1 - \sin x}\\
&= \frac{1 + \sin x}{2}\\
&\quad + \tan^2 x(1 + \sin x)\\
\frac{1 - \frac{1}{2}\sin^2 x}{1 - \sin x} &= (1 + \sin x)\left(\frac{1}{2} + \tan^2 x\right)\\
1 - \frac{1}{2}\sin^2 x &= (1 - \sin^2 x)\left(\frac{\cos^2 x + 2\sin^2 x}{2\cos^2 x}\right)\\
\frac{2 - \sin^2 x}{2} &= (1 - \sin^2 x)\frac{\cos^2 x + \sin^2 x + \sin^2 x}{2(1 - \sin^2 x)}\\
2 - \sin^2 x &= 1 + \sin^2 x
\end{aligned}
\]
$\sin^2 x = \dfrac{1}{2}$, $\sin x = \pm\dfrac{\sqrt{2}}{2}$, $x = \pm\dfrac{\pi}{4} + k\pi$, $(k \in \mathbb{Z})$ ជាចម្លើយរបស់សមីការ ។
\end{example}

\begin{example}[លំហាត់គំរូ ៣]
ដោះស្រាយសមីការ $3\sin^2 x + 2\sin x\cos x = 2$ ។
\end{example}
\ans
សមីការសរសេរ $3\sin^2 x + 2\sin x\cos x = 2(\sin^2 x + \cos^2 x)$
\[
\begin{gathered}
\sin^2 x + 2\sin x\cos x - 2\cos^2 x = 0\\
\text{ចែកអង្គទាំងពីរនៃសមីការនឹង } \cos^2 x \ne 0 \text{ ។}
\end{gathered}
\]
គេបាន $\tan^2 x + 2\tan x - 2 = 0$, $\tan x = -1 \pm \sqrt{3}$ ។ តាង $\tan\alpha = -1 \pm \sqrt{3}$ ៖ $\tan x = \tan\alpha$, $x = \alpha + k\pi$, $(k \in \mathbb{Z})$ ។

\begin{example}[លំហាត់គំរូ ៤]
រកតម្លៃ $m$ ដែលនាំឱ្យសមីការ $1 + m\cos x = m^2 - \cos^2 x$ មានឫស ។
\end{example}
\ans
តាង $\cos x = t$ ដែល $-1 \le t \le 1$ ។ គេបាន $t^2 + mt + 1 - m^2 = 0$ ។ តាង $f(t) = t^2 + mt + 1 - m^2$ ។

សមីការមានឫសមួយលើចន្លោះ $[-1, 1]$ កាលណា $f(-1) \times f(1) \le 0$ ។ គេបាន $(2 - m - m^2)(2 + m - m^2) \le 0$
\[
\begin{gathered}
(m^2 + m - 2)(m^2 - m - 2) \le 0,\\
-2 \le m \le -1 \ \text{ឬ} \ (1 \le m \le 2) \quad (1)
\end{gathered}
\]
សមីការមានឫសពីរលើ $[-1, 1]$ កាលណា
\[
\begin{cases} \Delta \ge 0\\ f(-1) > 0\\ f(1) > 0\\ -1 < \frac{S}{2} < 1 \end{cases}
\quad
\begin{cases} m^2 - 4(1 - m^2) \ge 0\\ -m^2 - m + 2 > 0\\ -m^2 + m + 2 > 0\\ -1 < -\frac{m}{2} < 1 \end{cases}
\quad
\begin{cases} m \le -\frac{2}{\sqrt{5}}, \ \frac{2}{\sqrt{5}} \le m\\ -2 < m < 1\\ -1 < m < 2\\ -2 < m < 2 \end{cases}
\]
\[ -1 < m \le -\frac{2\sqrt{5}}{5}, \quad \frac{2\sqrt{5}}{5} \le m < 1 \quad (2) \]
តាម (1) និង (2) នាំឱ្យសមីការមានឫសកាលណា $-2 \le m \le -\dfrac{2\sqrt{5}}{5}$ ឬ $\dfrac{2\sqrt{5}}{5} \le m \le 2$ ។

\begin{pro}
\begin{enumerate}[1]
  \item ដោះស្រាយសមីការខាងក្រោម ៖
  \begin{enumerate}[a]
    \item $7\cos^2 x + 6\sqrt{3}\sin x\cos x - \sin^2 x + 2 = 0$
    \item $2\cos^2 x + \cot^2 x = \dfrac{\sin^3 x + 1}{\sin^2 x}$
    \item $\tan^3 x + \tan^2 x - 3\tan x = 3$
    \item $\sin^2 x + 3\cos^2 x - 2\sin 2x = 0$ ។
  \end{enumerate}
  \item រកតម្លៃ $m$ ដើម្បីឱ្យសមីការ $\sin 4x = m\tan x$ មានឫស ។
\end{enumerate}
\ans
\begin{enumerate}[1]
  \item
  \begin{enumerate}[a]
    \item $\tan x = -3\sqrt{3}$ ៖ $x = \alpha + k\pi$ ដែល $\tan\alpha = -3\sqrt{3}$, $(k \in \mathbb{Z})$
    \item $x = \dfrac{\pi}{6} + 2k\pi$, $x = \dfrac{5\pi}{6} + 2k\pi$, $x = -\dfrac{\pi}{2} + 2k\pi$, $(k \in \mathbb{Z})$
    \item $x = \dfrac{3\pi}{4} + k\pi$, $x = \dfrac{\pi}{3} + k\pi$, $x = \dfrac{2\pi}{3} + k\pi$, $(k \in \mathbb{Z})$
    \item $x = \dfrac{\pi}{4} + k\pi$ ឬ $x = \alpha + k\pi$ ដែល $\tan\alpha = 3$, $(k \in \mathbb{Z})$ ។
  \end{enumerate}
  \item $-\dfrac{1}{2} \le m < 4$ \ ($x \ne k\pi$) ។
\end{enumerate}
\end{pro}

\section{ប្រព័ន្ធសមីការត្រីកោណមាត្រ}

\begin{example}[ឧទាហរណ៍ ១]
ដោះស្រាយប្រព័ន្ធសមីការ $\begin{cases} x + y = \dfrac{\pi}{3}\\ \sin x + \sin y = 1 \end{cases}$ ។

ប្រើរូបមន្តបំប្លែងពីផលបូកទៅជាផលគុណ គេបានសមីការ $\sin x + \sin y = 2\sin\left(\dfrac{x + y}{2}\right)\cos\left(\dfrac{x - y}{2}\right) = 1$ ដោយ $x + y = \dfrac{\pi}{3}$ គេបាន $2\sin\dfrac{\pi}{6}\cos\left(\dfrac{x - y}{2}\right) = 1$ ។

គេមាន $\sin\dfrac{\pi}{6} = \dfrac{1}{2}$ នាំឱ្យ $\cos\left(\dfrac{x - y}{2}\right) = 1 = \cos 0$ គេបាន $\dfrac{x - y}{2} = 2k\pi$ ឬ $x - y = 4k\pi$ ។
\[ \text{គេបានប្រព័ន្ធ} \quad \begin{cases} x + y = \frac{\pi}{3}\\ x - y = 4k\pi \end{cases}, \qquad \begin{cases} x = \frac{\pi}{6} + 2k\pi\\ y = \frac{\pi}{6} - 2k\pi \end{cases}, \ (k \in \mathbb{Z}) \text{ ។} \]
ដូច្នេះ ប្រព័ន្ធសមីការមានគូចម្លើយ $\left(x = \dfrac{\pi}{6} + 2k\pi, \ y = \dfrac{\pi}{6} - 2k\pi\right)$, $(k \in \mathbb{Z})$ ។
\end{example}

\begin{example}[ឧទាហរណ៍ ២]
ដោះស្រាយប្រព័ន្ធសមីការ $\begin{cases} x + y = \dfrac{\pi}{2} & (1)\\ \tan x + \tan y = \sqrt{3} + \dfrac{\sqrt{3}}{3} & (2) \end{cases}$

លក្ខខណ្ឌ $x \ne \dfrac{\pi}{2} + k\pi$, $y \ne \dfrac{\pi}{2} + k\pi$, $(k \in \mathbb{Z})$ ។ ដោយ $\tan y = \cot x$ តាមមុំបំពេញ គេបាន
\[ \begin{cases} x + y = \frac{\pi}{2}\\ \tan^2 x - \left(\sqrt{3} + \frac{\sqrt{3}}{3}\right)\tan x + 1 = 0 \end{cases} \]
សមីការ $\tan^2 x - \left(\sqrt{3} + \dfrac{\sqrt{3}}{3}\right)\tan x + 1 = 0$ មានឫសពីរគឺ $t_1 = \sqrt{3}$, $t_2 = \dfrac{\sqrt{3}}{3}$ ។ ដូច្នេះ គេបាន
\[ \begin{cases} x + y = \frac{\pi}{2}\\ \tan x = \sqrt{3} = \tan\frac{\pi}{3} \end{cases} \ \text{ឬ} \ \begin{cases} x + y = \frac{\pi}{2}\\ \tan x = \frac{\sqrt{3}}{3} = \tan\frac{\pi}{6} \end{cases} \]
\[ \begin{cases} x = \frac{\pi}{3} + k\pi\\ y = \frac{\pi}{6} - k\pi \end{cases} \ \text{ឬ} \ \begin{cases} x = \frac{\pi}{6} + k\pi\\ y = \frac{\pi}{3} - k\pi \end{cases}, \ (k \in \mathbb{Z}) \text{ ជាចម្លើយរបស់ប្រព័ន្ធសមីការ ។} \]
\end{example}

\begin{example}[លំហាត់គំរូ ៥]
រកតម្លៃ $m$ ដែលនាំឱ្យប្រព័ន្ធសមីការខាងក្រោមមានឫស
\[ \begin{cases} x - y = m & (1)\\ 2(\cos 2x + \cos 2y) - 1 - 4\cos^2 m = 0 & (2) \end{cases} \]
\end{example}
\ans
គេបានសមីការ (2) សរសេរ ៖ $4\cos(x + y)\cos(y - x) - 1 - 4\cos^2 m = 0$ តាម (1) ជំនួសចូល គេបាន ៖
\[
\begin{gathered}
4\cos^2 m - 4\cos(x + y)\cos m + 1 = 0,\\
[2\cos m - \cos(x + y)]^2 + \sin^2(x + y) = 0
\end{gathered}
\]
\[ \begin{cases} \sin(x + y) = 0\\ \cos(x + y) = 2\cos m \end{cases}, \quad \begin{cases} x + y = 2k\pi\\ \cos m = \frac{1}{2} \end{cases} \ \text{ឬ} \ \begin{cases} x + y = \pi + 2k\pi\\ \cos m = -\frac{1}{2} \end{cases} \]
\begin{itemize}
  \item ចំពោះ $\cos m = \dfrac{1}{2}$, $m = \pm\dfrac{\pi}{3} + 2k\pi$ នោះ $\cos(x + y) = 1$ គេបាន $\begin{cases} x + y = 2k\pi\\ x - y = m \end{cases}$ ។
  \item ចំពោះ $\cos m = -\dfrac{1}{2}$, $m = \pm\dfrac{2\pi}{3} + 2k\pi$ នោះ $\cos(x + y) = -1$ គេបាន $\begin{cases} x + y = (2k + 1)\pi\\ x - y = m \end{cases}$, \ $\begin{cases} x = \frac{m}{2} + (2k + 1)\frac{\pi}{2}\\ y = -\frac{m}{2} + (2k + 1)\frac{\pi}{2} \end{cases}$, $(k \in \mathbb{Z})$ ។
\end{itemize}
ដូច្នេះ ៖
\begin{itemize}
  \item $m = \pm\dfrac{\pi}{3} + 2k\pi$ ប្រព័ន្ធមានឫស $\begin{cases} x = \frac{m}{2} + k\pi\\ y = -\frac{m}{2} + k\pi \end{cases}$, $(k \in \mathbb{Z})$
  \item $m = \pm\dfrac{2\pi}{3} + 2k\pi$ ប្រព័ន្ធមានឫស $\begin{cases} x = \frac{m}{2} + (2k + 1)\frac{\pi}{2}\\ y = -\frac{m}{2} + (2k + 1)\frac{\pi}{2} \end{cases}$, $(k \in \mathbb{Z})$ ។
\end{itemize}

\begin{pro}
\begin{enumerate}[1]
  \item ដោះស្រាយប្រព័ន្ធសមីការ
  \begin{tasks}(2)
    \task $\begin{cases} \cos^2 x + \cos^2 y = \frac{1}{4}\\ x + y = \frac{5\pi}{6} \end{cases}$
    \task $\begin{cases} \sin x\cos y = \frac{1}{4}\\ 3\tan x = \tan y \end{cases}$ ។
  \end{tasks}
  \item គេឱ្យប្រព័ន្ធសមីការ $\begin{cases} \sin x\cos 2y = m^2 + 1\\ \cos x\sin 2y = m \end{cases}$ ។
  \begin{enumerate}[a]
    \item បញ្ជាក់តម្លៃ $m$ ដែលនាំឱ្យប្រព័ន្ធសមីការមានឫស ដែល $m$ ជាចំនួនគត់វិជ្ជមាន ។
    \item ដោះស្រាយប្រព័ន្ធសមីការតាមតម្លៃ $m$ ដែលរកបាននោះ ។
  \end{enumerate}
\end{enumerate}
\ans
\begin{enumerate}[1]
  \item ក. $x - y = \pm\dfrac{\pi}{6} + 2k\pi$ ៖ $\left(x = \dfrac{\pi}{2} + k\pi, \ y = \dfrac{\pi}{3} - k\pi\right)$ ឬ $\left(x = \dfrac{\pi}{3} + k\pi, \ y = \dfrac{\pi}{2} - k\pi\right)$, $(k \in \mathbb{Z})$ ។

  ខ. $\sin(x + y) = 1$, $\sin(x - y) = -\dfrac{1}{2}$ ៖ $\left(x = \dfrac{\pi}{6} + n\pi, \ y = \dfrac{\pi}{3} + (2k - n)\pi\right)$ ឬ $\left(x = \dfrac{5\pi}{6} + n\pi, \ y = -\dfrac{\pi}{3} + (2k - n)\pi\right)$, $(k, n \in \mathbb{Z})$ ។
  \item ក. $m = 0$ \qquad ខ. $\left(x = \dfrac{\pi}{2}, \ y = 0\right)$ (តាម $2\pi$) ។
\end{enumerate}
\end{pro}

\section{វិសមីការត្រីកោណមាត្រ}

\subsection{វិសមីការដែលមានរាង \texorpdfstring{$a\cos x + b\sin x > c$}{a cos x + b sin x > c}}

\begin{example}[ឧទាហរណ៍ ១]
ដោះស្រាយវិសមីការ $\sin x - \cos x > 0$ ។
\[
\begin{gathered}
\sqrt{2}\left(\frac{\sqrt{2}}{2}\sin x - \frac{\sqrt{2}}{2}\cos x\right) > 0,\\
\sqrt{2}\left(\sin x\cos\frac{\pi}{4} - \sin\frac{\pi}{4}\cos x\right) > 0,\\
\sqrt{2}\sin\left(x - \frac{\pi}{4}\right) > 0 \text{ ។}
\end{gathered}
\]
$\sin\left(x - \dfrac{\pi}{4}\right) > 0$ ត្រូវនឹងចុងធ្នូនៅលើកន្លះរង្វង់ខាងលើ ។

ដូច្នេះ វិសមីការមានសំណុំចម្លើយ $2k\pi < x - \dfrac{\pi}{4} < \pi + 2k\pi$ ឬ $\dfrac{\pi}{4} + 2k\pi < x < \dfrac{5\pi}{4} + 2k\pi$, $(k \in \mathbb{Z})$ ។
\end{example}

\begin{note}[សម្គាល់]
បើវិសមីការដែលមានទម្រង់ជាកន្សោមប្រភាគ ឬវិសមីការដឺក្រេទី $2$ គេត្រូវដោះស្រាយរកសំណុំចម្លើយដោយសិក្សាសញ្ញា ។
\end{note}

\begin{example}[ឧទាហរណ៍ ២]
ដោះស្រាយវិសមីការ $\sin x + \dfrac{1}{\sin x} \ge \dfrac{5}{2}$ ។

គេមាន $\sin x + \dfrac{1}{\sin x} - \dfrac{5}{2} \ge 0$
\[
\begin{gathered}
\frac{2\sin^2 x - 5\sin x + 2}{2\sin x} \ge 0,\\
\frac{(\sin x - 2)(2\sin x - 1)}{2\sin x} \ge 0
\end{gathered}
\]
ដោយ $\sin x - 2 < 0$ ចំពោះគ្រប់ $x \in \mathbb{R}$ គេបាន $\dfrac{2\sin x - 1}{2\sin x} \le 0$ \ (1) ។ វិសមីការមានចម្លើយលុះត្រាតែ $\sin x \ne 0$ ឬ $x \ne k\pi$ ។

តាង $\sin x = X$ ដែល $-1 \le X \le 1$ វិសមីការ (1) ទៅជា $\dfrac{2X - 1}{2X} \le 0$ ។ សិក្សាសញ្ញាគេបាន
\[
\begin{array}{|c|ccccccc|}\hline
X & -1 & & 0 & & \frac{1}{2} & & 1\\ \hline
2X - 1 & & - & & - & 0 & + & \\ \hline
X & & - & 0 & + & & + & \\ \hline
\frac{2X - 1}{X} & & + & \| & - & 0 & + & \\ \hline
\end{array}
\]
គេបាន $0 < X \le \dfrac{1}{2}$, \ $0 < \sin x \le \dfrac{1}{2}$ ត្រូវនឹង $2k\pi < x \le \dfrac{\pi}{6} + 2k\pi$ ឬ $\dfrac{5\pi}{6} + 2k\pi \le x < \pi + 2k\pi$, $(k \in \mathbb{Z})$ ។
\end{example}

\begin{example}[លំហាត់គំរូ ៦]
ដោះស្រាយវិសមីការ $2\cos 2x + \sin^2 x\cos x + \sin x\cos^2 x > 2(\sin x + \cos x)$ ។
\end{example}
\ans
គេមាន
\[
\begin{aligned}
&2(\cos^2 x - \sin^2 x) + \sin x\cos x(\sin x + \cos x)\\
&\quad - 2(\sin x + \cos x) > 0\\
&(\sin x + \cos x)[2(\cos x - \sin x) + \sin x\cos x - 2] > 0
\end{aligned}
\]
តាង $f(x) = (\sin x + \cos x)[2(\cos x - \sin x) + \sin x\cos x - 2]$ ។ ដោយ $f(x)$ ជាអនុគមន៍ដែលមានខួប $2\pi$ នោះគេអាចសិក្សាវិសមីការក្នុងចន្លោះ $(0, 2\pi)$ ។

គេបាន $f(x) = 0$, \ $\cos x + \sin x = 0$ \ (1) ឬ $2(\cos x - \sin x) + \sin x\cos x - 2 = 0$ \ (2)

(1) \ $\tan x = -1$, \ $\left(x = \dfrac{3\pi}{4}, \ x = \dfrac{7\pi}{4}\right)$ ។

(2) \ តាង $t = \cos x - \sin x = \sqrt{2}\cos\left(x + \dfrac{\pi}{4}\right)$ ដែល $-\sqrt{2} \le t \le \sqrt{2}$ ។ $t^2 = 1 - 2\sin x\cos x$, \ $\sin x\cos x = \dfrac{1 - t^2}{2}$ ។

(2) អាចសរសេរ $2t + \dfrac{1 - t^2}{2} - 2 = 0$ ឬ $t^2 - 4t + 3 = 0$ ។ ដោះស្រាយសមីការនេះ គេបាន $t = 1$, $t = 3$ (មិនយក)

$t = 1$, \ $\sqrt{2}\cos\left(x + \dfrac{\pi}{4}\right) = 1$, \ $\cos\left(x + \dfrac{\pi}{4}\right) = \dfrac{\sqrt{2}}{2} = \cos\dfrac{\pi}{4}$ ។ គេបាន $x = 0$, $x = \dfrac{3\pi}{2}$ ។

គេបានសញ្ញាដូចខាងក្រោម ៖
\[
\begin{array}{|c|ccccccccc|}\hline
x & 0 & & \frac{3\pi}{4} & & \frac{3\pi}{2} & & \frac{7\pi}{4} & & 2\pi\\ \hline
f(x) & & - & 0 & + & 0 & - & 0 & + & \\ \hline
\end{array}
\]
ដូច្នេះ វិសមីការមានសំណុំចម្លើយ $\dfrac{3\pi}{4} + 2k\pi < x < \dfrac{3\pi}{2} + 2k\pi$ ឬ $\dfrac{7\pi}{4} + 2k\pi < x < 2\pi + 2k\pi$, $(k \in \mathbb{Z})$ ។

\begin{example}[លំហាត់គំរូ ៧]
បង្ហាញថាក្នុងត្រីកោណ $ABC$ គេបាន $\cos A + \cos B + \cos C \le \dfrac{3}{2}$ ។
\end{example}
\ans
គេមាន
\[
\begin{aligned}
&\frac{3}{2} - (\cos A + \cos B + \cos C)\\
&= \frac{1}{2}[3 - 2(\cos A + \cos B + \cos C)]\\
&= \frac{1}{2}\left[1 - 4\cos\frac{A + B}{2}\cos\frac{A - B}{2} + 2(1 - \cos C)\right]\\
&= \frac{1}{2}\left[1 - 4\cos\frac{A + B}{2}\cos\frac{A - B}{2} + 4\sin^2\frac{C}{2}\right]\\
&= \frac{1}{2}\left[1 - 4\sin\frac{C}{2}\cos\frac{A - B}{2} + 4\sin^2\frac{C}{2}\right]\\
&\quad \left(\text{ព្រោះ } \cos\frac{A + B}{2} = \sin\frac{C}{2}\right)\\
&= \frac{1}{2}\left[\left(2\sin\frac{C}{2} - \cos\frac{A - B}{2}\right)^2 + 1 - \cos^2\frac{A - B}{2}\right]\\
&= \frac{1}{2}\left[\left(2\sin\frac{C}{2} - \cos\frac{A - B}{2}\right)^2 + \sin^2\frac{A - B}{2}\right]\\
&\ge 0
\end{aligned}
\]
ដូច្នេះ $\dfrac{3}{2} \ge \cos A + \cos B + \cos C$ ឬ $\cos A + \cos B + \cos C \le \dfrac{3}{2}$ ។

\begin{pro}
\begin{enumerate}[1]
  \item ដោះស្រាយវិសមីការខាងក្រោម ៖
  \begin{tasks}(2)
    \task $\sin x + \cos x > \cos\dfrac{\pi}{4}$
    \task $\dfrac{2\cos x - 5\sin x}{\cos x} > 0$
    \task $\dfrac{\cos x}{1 - 3\cos x} < \dfrac{1 - \cos x}{1 - 9\cos^2 x}$
    \task $2\sin^2\left(x + \dfrac{\pi}{4}\right) + \sqrt{3}\cos 2x > 0$ ។
  \end{tasks}
  \item បង្ហាញថាក្នុងត្រីកោណ $ABC$ គេបាន $\sin^2 A + \sin^2 B + \sin^2 C \le \dfrac{9}{4}$ ។
\end{enumerate}
\ans
\begin{enumerate}[a]
  \item $-\dfrac{\pi}{12} + 2k\pi < x < \dfrac{7\pi}{12} + 2k\pi$, $(k \in \mathbb{Z})$
  \item $-\dfrac{\pi}{2} + k\pi < x < \alpha + k\pi$ ដែល $\tan\alpha = \dfrac{2}{5}$, $(k \in \mathbb{Z})$
  \item $-\alpha + 2k\pi < x < \alpha + 2k\pi$ ដែល $\cos\alpha = -\dfrac{1}{3}$ ($\cos x \ne \pm\dfrac{1}{3}$), $(k \in \mathbb{Z})$
  \item $-\dfrac{\pi}{4} + k\pi < x < \dfrac{5\pi}{12} + k\pi$, $(k \in \mathbb{Z})$ ។
\end{enumerate}
\end{pro}

\section*{មេរៀនសង្ខេប ជំពូកទី ៣}
\begin{note}[មេរៀនសង្ខេប]
\begin{enumerate}[1]
  \item \textbf{រូបមន្តផលបូកនិងផលដក} \
  $\cos(\alpha \pm \beta) = \cos\alpha\cos\beta \mp \sin\alpha\sin\beta$, \ $\sin(\alpha \pm \beta) = \sin\alpha\cos\beta \pm \sin\beta\cos\alpha$, \ $\tan(\alpha \pm \beta) = \dfrac{\tan\alpha \pm \tan\beta}{1 \mp \tan\alpha\tan\beta}$ ។
  \item \textbf{រូបមន្តមុំឌុប} \ $\sin 2\alpha = 2\sin\alpha\cos\alpha$, \ $\cos 2\alpha = \cos^2\alpha - \sin^2\alpha = 1 - 2\sin^2\alpha = 2\cos^2\alpha - 1$ ។
  \item \textbf{រូបមន្តកន្លះមុំ} \ $\sin\alpha = 2\sin\dfrac{\alpha}{2}\cos\dfrac{\alpha}{2}$, \ $\cos^2\dfrac{\alpha}{2} = \dfrac{1 + \cos\alpha}{2}$, \ $\sin^2\dfrac{\alpha}{2} = \dfrac{1 - \cos\alpha}{2}$, \ $\tan\dfrac{\alpha}{2} = \dfrac{\sin\alpha}{1 + \cos\alpha}$ ។ បើ $\tan\dfrac{\alpha}{2} = t$ គេបាន $\cos\alpha = \dfrac{1 - t^2}{1 + t^2}$, $\sin\alpha = \dfrac{2t}{1 + t^2}$, $\tan\alpha = \dfrac{2t}{1 - t^2}$ ។
  \item \textbf{រូបមន្តបំប្លែង} \ $\sin\alpha\sin\beta = \dfrac{1}{2}[\cos(\alpha - \beta) - \cos(\alpha + \beta)]$, \ $\cos\alpha\cos\beta = \dfrac{1}{2}[\cos(\alpha + \beta) + \cos(\alpha - \beta)]$, \ $\sin\alpha\cos\beta = \dfrac{1}{2}[\sin(\alpha + \beta) + \sin(\alpha - \beta)]$ ។
  \item \textbf{បំប្លែងពីផលបូកទៅផលគុណ} \ $\cos p + \cos q = 2\cos\dfrac{p + q}{2}\cos\dfrac{p - q}{2}$, \ $\cos p - \cos q = -2\sin\dfrac{p + q}{2}\sin\dfrac{p - q}{2}$, \ $\sin p + \sin q = 2\sin\dfrac{p + q}{2}\cos\dfrac{p - q}{2}$, \ $\sin p - \sin q = 2\sin\dfrac{p - q}{2}\cos\dfrac{p + q}{2}$, \ $\tan p \pm \tan q = \dfrac{\sin(p \pm q)}{\cos p\cos q}$ ។
  \item \textbf{អនុគមន៍ច្រាសនៃអនុគមន៍ត្រីកោណមាត្រ} \ បើ $a = f(b)$ នោះអនុគមន៍ច្រាសគឺ $b = f^{-1}(a)$ ។ $\sin^{-1}(-x) = -\sin^{-1}(x)$, $\cos^{-1}(-x) = \pi - \cos^{-1}(x)$, $\tan^{-1}(-x) = -\tan^{-1}(x)$ ។ $\sin^{-1}(\sin x) = x$, $x \in \left[-\dfrac{\pi}{2}, \dfrac{\pi}{2}\right]$; $\cos^{-1}(\cos x) = x$, $x \in [0, \pi]$; $\sin(\sin^{-1}x) = x$, $\cos(\cos^{-1}x) = x$, $x \in [-1, 1]$ ។
  \item សមីការ $\cos x = \cos\alpha$ មានចម្លើយពីរគឺ $x = \alpha + 2k\pi$ និង $x = -\alpha + 2k\pi$ $(k \in \mathbb{Z})$ ។
  \item សមីការ $\sin x = \sin\alpha$ មានចម្លើយពីរគឺ $x = \alpha + 2k\pi$ និង $x = \pi - \alpha + 2k\pi$, $(k \in \mathbb{Z})$ ។
  \item សមីការ $\tan x = \tan\alpha$ មានចម្លើយ $x = \alpha + k\pi$; សមីការ $\cot x = \cot\alpha$ មានចម្លើយ $x = \alpha + k\pi$, $(k \in \mathbb{Z})$ ។
  \item វិសមីការ ៖ $\sin x \le a$ ៖ $-\pi - \alpha + 2k\pi \le x \le \alpha + 2k\pi$; \ $\sin x \ge a$ ៖ $\alpha + 2k\pi \le x \le \pi - \alpha + 2k\pi$; \ $\cos x \le a$ ៖ $\alpha + 2k\pi \le x \le 2\pi - \alpha + 2k\pi$; \ $\cos x \ge a$ ៖ $-\alpha + 2k\pi \le x \le \alpha + 2k\pi$; \ $\tan x \le a$ ៖ $-\dfrac{\pi}{2} + k\pi < x \le \alpha + k\pi$; \ $\tan x \ge a$ ៖ $\alpha + k\pi \le x < \dfrac{\pi}{2} + k\pi$ ។
  \item ដើម្បីដោះស្រាយសមីការ $a\cos x + b\sin x = c$ គេបំប្លែងជារាង $\cos(x - \theta) = \dfrac{c}{r}$ ដែល $r = \sqrt{a^2 + b^2}$, $\cos\theta = \dfrac{a}{r}$, $\sin\theta = \dfrac{b}{r}$ ឬជំនួស $t = \tan\dfrac{x}{2}$ $(x \ne \pi + 2k\pi)$ ទៅជាសមីការដឺក្រេទី $2$ ៖ $(c + a)t^2 - 2bt + c - a = 0$ ។
  \item ដើម្បីដោះស្រាយសមីការ $a\cos^2 x + b\cos x + c = 0$, $a\sin^2 x + b\sin x + c = 0$, $a\tan^2 x + b\tan x + c = 0$ \ldots គេតាងអនុគមន៍ត្រីកោណមាត្រដោយ $t$, ដាក់លក្ខខណ្ឌ, ដោះស្រាយសមីការដឺក្រេទី $2$ រួចយកឫសមកដោះស្រាយតាមរូបមន្តសមីការត្រីកោណមាត្រ ។
  \item ដើម្បីដោះស្រាយវិសមីការ $a\cos x + b\sin x < c$ \ldots គេបំប្លែងអង្គទី $1$ ជារាង $\cos(x - \theta)$ ឬ $\sin(x - \theta)$ ។ បើវិសមីការមានទម្រង់ជាកន្សោមប្រភាគ ឬដឺក្រេទី $2$ គេដោះស្រាយដោយសិក្សាសញ្ញា ។
\end{enumerate}
\end{note}

\section*{លំហាត់ជំពូកទី ៣}
\begin{enumerate}[1]
  \item អង្កត់ធ្នូ $AB$ ចែករង្វង់មួយដែលមានកាំស្មើនឹង $2$ cm ជាពីរអង្កត់ ។ គេដឹងថា $AB$ ស្កាត់ធ្នូដែលមានមុំផ្ចិត $= 60^\circ$ ។ គណនា
  \begin{enumerate}[a]
    \item ក្រឡាផ្ទៃចម្រៀកថាស $AOB$ ។
    \item ក្រឡាផ្ទៃត្រីកោណ $AOB$ ។
    \item ក្រឡាផ្ទៃអង្កត់ថាសតូច (ផ្នែកតូច) ។
  \end{enumerate}
  \item គណនា
  \begin{tasks}(3)
    \task $\sin 960^\circ$
    \task $\cos 1215^\circ$
    \task $\sin 2940^\circ$
    \task $\cos\dfrac{13\pi}{6}$
    \task $\tan\left(-\dfrac{19}{4}\pi\right)$
    \task $\sin\left(-\dfrac{17}{3}\pi\right)$ ។
  \end{tasks}
  \item គណនា
  \begin{enumerate}[a]
    \item $\sin\theta + \sin\left(\theta + \dfrac{\pi}{2}\right) + \sin(\theta + \pi) + \sin\left(\theta + \dfrac{3}{2}\pi\right)$
    \item $\cos\dfrac{20}{3}\pi\tan\dfrac{7}{4}\pi + \sin\dfrac{5}{2}\pi\cos(-3\pi)$
    \item $\sin\left(\dfrac{\pi}{2} - \theta\right) + \cos\left(\dfrac{\pi}{2} + 2\theta\right) + \sin(\pi - 2\theta) + \cos(\pi - \theta)$ ។
  \end{enumerate}
  \item គេមាន $\tan\alpha + \tan\beta + \tan\gamma = \tan\alpha\tan\beta\tan\gamma$ $\left(-\dfrac{\pi}{2} < \alpha, \beta, \gamma < \dfrac{\pi}{2}\right)$ ។ គណនាមុំ $(\alpha + \beta + \gamma)$ ។
  \item គេមាន $\alpha = \cos 10^\circ$, $\beta = \cos 50^\circ$, $\gamma = \cos 70^\circ$ ។ គណនា $\alpha - \beta - \gamma$, $\alpha^2 - \beta\gamma$ និង $\alpha^2 + \beta^2 + \gamma^2$ ។
  \item បង្ហាញថា
  \begin{tasks}(2)
    \task $(1 - \sin A + \cos A)^2 = 2(1 - \sin A)(1 + \cos A)$
    \task $\cot A + \tan A = \dfrac{1}{\sin A\cos A}$
    \task $(\tan\theta + \sin\theta)(\tan\theta - \sin\theta) = \tan^2\theta\sin^2\theta$
    \task $\left(1 + \dfrac{1}{\sin\theta}\right)(1 - \sin\theta) = \cos\theta\cot\theta$ ។
  \end{tasks}
  \item ដោះស្រាយសមីការខាងក្រោម ៖
  \begin{tasks}(2)
    \task $4\cos^2 x - \cos 2x - 4\cos x + 1 = 0$
    \task $\sin^2 2x - \cos^2 8x = 0.5\cos 10x$
    \task $1 + \sin x + \cos x + \sin 2x + \cos 2x = 0$
    \task $\sin^4 x + \cos^4 x - \cos 2x + \dfrac{1}{4}\sin^2 2x - 2 = 0$
    \task $\sin^{2007} x + \cos^{2007} x = 1$
    \task $\tan x = \cot x + 2\cot^2 2x$ ។
  \end{tasks}
  \item
  \begin{enumerate}[a]
    \item គេឱ្យ $2\sin A\cos A + (\cos A + \sin A)^2 - (2\cos A + \sin A)^2 = p\sin^2 A + q$ ។ រកតម្លៃនៃចំនួនថេរ $p$ និង $q$ ។
    \item គេឱ្យ $\sin x\cos x(5\tan x + 2\cot x) = a + b\sin^2 x$ ។ រកតម្លៃនៃចំនួនថេរ $a$ និង $b$ ។
  \end{enumerate}
  \item ដោះស្រាយវិសមីការខាងក្រោម ដោយដឹងថា $0 \le x < 2\pi$
  \begin{tasks}(2)
    \task $\cos\left(x + \dfrac{\pi}{3}\right) \ge -\dfrac{1}{2}$
    \task $\tan\left(x - \dfrac{\pi}{6}\right) \ge -1$ ។
  \end{tasks}
  \item ដោះស្រាយសមីការនិងវិសមីការខាងក្រោម ៖
  \begin{tasks}(2)
    \task $2\sin^2\theta - \sqrt{2}\cos\theta = 0$
    \task $2\cos^2\theta + \sqrt{3}\sin\theta + 1 > 0$ ។
  \end{tasks}
  \item គេឱ្យ $\triangle ABC$ មានមុំ $3$ ជាមុំស្រួច ។
  \begin{enumerate}[a]
    \item បង្ហាញថា $\tan A + \tan B + \tan C = \tan A \cdot \tan B \cdot \tan C$
    \item តាង $T = \tan A + \tan B + \tan C$ ។ បង្ហាញថា $T \ge 3\sqrt{3}$ ។
  \end{enumerate}
  \item ដោះស្រាយសមីការខាងក្រោម ៖
  \begin{tasks}(2)
    \task $2\cos 3x + \sqrt{3}\sin x + \cos x = 0$
    \task $\sin^6 x + \cos^3 x = 2(\sin^6 x + \cos^8 x)$
    \task $\dfrac{\sin 5x}{5\sin x} = 1$
    \task $\sin^2 x + \sin^2 3x = \cos^2 2x + \cos^2 4x$
    \task $\cos^3 x + \dfrac{3\sqrt{2}}{4}\sin 2x - 2\cos x = 0$
    \task $\sin x + \sin 2x + \sin 3x + \sin 4x + \sin 5x + \sin 6x = 0$ ។
  \end{tasks}
  \item ដោះស្រាយវិសមីការខាងក្រោម ៖
  \begin{enumerate}[a]
    \item $2\cos^2 x - \cos x + 1 \le 0$ ក្នុង $[0, \pi]$
    \item $\dfrac{2\sin^2 x - \sin x - 1}{\sin x} > 0$ ក្នុង $[0, \pi]$
    \item $\dfrac{\sin x - \cos x + 1}{\sin x + \cos x - 1} > 0$ ។
  \end{enumerate}
  \item
  \begin{enumerate}[a]
    \item ពន្លាតកន្សោម $\left(x + \dfrac{1}{2}\right)(x - 8)(x - 1)$ ។
    \item ដោះស្រាយសមីការ $2\sin^3 x - 17\sin^2 x + 7\sin x + 8 = 0$ ។
  \end{enumerate}
  \item គេឱ្យសមីការ $\cos 2x - (2m + 1)\cos x + m + 1 = 0$ \ (1) ។
  \begin{enumerate}[a]
    \item ដោះស្រាយសមីការ (1) កាលណា $m = 1$
    \item រកតម្លៃ $m$ ដែលធ្វើឱ្យសមីការមានឫសនៅចន្លោះ $[0, \pi]$ ។
  \end{enumerate}
  \item គេមានវិសមីការ $(E)$ ៖ $2\sin^2 x - 5\sin x + 2 > 0$ ។ គេតាង $X = \sin x$ វិសមីការសរសេរ $2X^2 - 5X + 2 > 0$ ។
  \begin{enumerate}[a]
    \item ដាក់ជាផលគុណកត្តាកន្សោម $2x^2 - 5x + 2$ ។
    \item បង្ហាញថាវិសមីការ $(E)$ សរសេរជារាង $2(\sin x - 2)\left(\sin x - \dfrac{1}{2}\right) > 0$ ។
    \item សិក្សាសញ្ញា $(\sin x - 2)\left(\sin x - \dfrac{1}{2}\right)$ នៅលើចន្លោះ $[0, 2\pi]$ ។
    \item រកសំណុំចម្លើយនៃវិសមីការ $(E)$ ។
  \end{enumerate}
  \item ដោះស្រាយប្រព័ន្ធសមីការត្រីកោណមាត្រខាងក្រោម ៖
  \begin{tasks}(3)
    \task $\begin{cases} \sin x + \sin y = \frac{3}{2}\\ \sin^2 x + \sin^2 y = \frac{5}{4} \end{cases}$
    \task $\begin{cases} \cos^3 x + \cos x + \sin y = 0\\ \sin^3 x - \sin y + \cos x = 0 \end{cases}$
    \task $\begin{cases} \frac{1 - \tan x}{1 + \tan x} = \tan y\\ x - y = \frac{\pi}{6} \end{cases}$ ។
  \end{tasks}
  \item គេមាន $\triangle ABC$ ដែលមានមុំនិងជ្រុងបំពេញលក្ខខណ្ឌ ៖ $\dfrac{1 + \cos A}{1 + \cos B} = \dfrac{2a + c}{2a - c}$ ។ បង្ហាញ $\triangle ABC$ ជាត្រីកោណសមបាត ។
  \item មុំ $A$, $B$, $C$ នៃ $\triangle ABC$ មួយមាន $\dfrac{A + C}{2} = B$ ។ រករង្វាស់មុំនៃត្រីកោណនោះបើគេដឹងថា $\sin A + \sin B + \sin C = \dfrac{3 + \sqrt{3}}{2}$ ។
\end{enumerate}

\section*{លំហាត់អនុវត្តន៍}
\begin{enumerate}[1]
  \item ដោះស្រាយសមីការខាងក្រោម ៖
  \begin{tasks}(2)
    \task $\cos x + \sqrt{3}\sin x = \cos 3x$
    \task $\sin 3x + 2\cos x - 2 = 0$
    \task $\sin 2x + \tan x = 2$
    \task $\sin 5x + \cos 5x = \sqrt{2}\cos 13x$
    \task $6\sin x - 2\cos^3 x = 5\sin 2x\cos x$
    \task $\sqrt{5\cos x - \cos 2x} + 2\sin x = 0$ ។
  \end{tasks}
  \item ដោះស្រាយសមីការខាងក្រោម ៖
  \begin{tasks}(2)
    \task $1 + \cot 2x = \dfrac{1 - \cos 2x}{\sin^2 2x}$
    \task $\cos^6 x + \sin^6 x = \dfrac{7}{16}$
    \task $(1 - \tan x)(1 + \sin 2x) = 1 + \tan x$
    \task $3\sin 3x - \sqrt{3}\cos 9x = 1 + 4\sin^3 3x$
    \task $1 + 3\cos x + \cos 2x = \cos 3x + 2\sin x\sin 2x$
    \task $\cos^4 x + \sin^6 x = \cos 2x$ ។
  \end{tasks}
  \item គេឱ្យសមីការ $\cos 2x - (2m + 1)\cos x + m + 1 = 0$ \ (1) ។
  \begin{enumerate}[a]
    \item ដោះស្រាយសមីការ (1) កាលណា $m = \dfrac{3}{2}$ ។
    \item រកតម្លៃ $m$ ដែលធ្វើឱ្យសមីការមានឫស $x$ នៅចន្លោះ $\dfrac{\pi}{2} < x < \dfrac{3\pi}{2}$ ។
  \end{enumerate}
  \item ដោះស្រាយប្រព័ន្ធសមីការត្រីកោណមាត្រខាងក្រោម ៖
  \begin{tasks}(3)
    \task $\begin{cases} \sin(x + y) = \frac{1}{2}\\ \cos(x - y) = \frac{\sqrt{2}}{2} \end{cases}$
    \task $\begin{cases} \sin x + \sin y = \sqrt{2}\\ \cos x + \sin y = \sqrt{2} \end{cases}$
    \task $\begin{cases} \sin(x + y) = \cos(x - y)\\ \tan x - \tan y = 1 \end{cases}$ ។
  \end{tasks}
  \item ដោះស្រាយវិសមីការខាងក្រោម ៖
  \begin{tasks}(2)
    \task $\sin^2\left(\dfrac{x}{2} - \dfrac{\pi}{4}\right) < \cos^2\dfrac{x}{2}$
    \task $6\sin^2 x - \sin x\cos x - \cos^2 x > 2$
    \task $\dfrac{\cos x}{1 + 2\cos x} > \dfrac{1 - \cos x}{1 - 2\cos x}$
    \task $\dfrac{1 - \sin x}{1 - 3\sin x} < \dfrac{1 + \sin x}{1 - 9\sin^2 x}$ ។
  \end{tasks}
  \item បង្ហាញថា \quad ក. $\sin^4\alpha + \cos^4\alpha \ge \dfrac{1}{2}$ \quad ខ. $\sin^6\alpha + \cos^6\alpha \ge \dfrac{1}{4}$ \quad គ. $\sin^8\alpha + \cos^8\alpha \ge \dfrac{1}{8}$ ។
  \item បង្ហាញថាក្នុង $\triangle ABC$ គេបាន ៖ $\sin\dfrac{A}{2} \cdot \sin\dfrac{B}{2} \cdot \sin\dfrac{C}{2} \le \dfrac{1}{8}$ ។
  \item បង្ហាញថាក្នុង $\triangle ABC$ ដែលបំពេញលក្ខខណ្ឌ $\tan A\tan B\tan^2\dfrac{C}{2} = 1$ ជាត្រីកោណកែងសមបាត ។
  \item
  \begin{enumerate}[a]
    \item នៅពេលដែល $\sin\theta - \cos\theta = \dfrac{1}{3}$ ហើយ $0^\circ \le \theta \le 180^\circ$ រកតម្លៃ $\sin\theta + \cos\theta$
    \item នៅពេលដែល $\sin\alpha + \sin\beta = \sqrt{3}$, $\cos\alpha + \cos\beta = 1$ ហើយ $0^\circ < \alpha \le \beta < 180^\circ$ រកតម្លៃ $\alpha$ និង $\beta$
    \item នៅពេលដែល $\sin\theta + \cos\theta = \dfrac{17}{13}$ ហើយ $45^\circ \le \theta \le 90^\circ$ រកតម្លៃ $\sin\theta$, $\cos\theta$ និង $\tan\theta$ ។
  \end{enumerate}
  \item ដោះស្រាយវិសមីការ បើ $-180^\circ \le x \le 180^\circ$ និង $\sin x + \sqrt{3}\cos x \ge 1$ ។
  \item ដោះស្រាយប្រព័ន្ធសមីការ $\begin{cases} \sin 2x + \cos y = 1\\ \sin y + \cos 2x = 1 \end{cases}$ \ $(0^\circ \le x < 360^\circ, \ 0^\circ \le y < 360^\circ)$ ។
  \item ដោះស្រាយប្រព័ន្ធសមីការ $\begin{cases} \sin x = \sqrt{2}\sin y\\ \tan x = \sqrt{2}\tan y \end{cases}$ \ $(0^\circ < x < 180^\circ, \ 0^\circ < y < 180^\circ)$ ។
  \item ឧបមាថាកន្លះបន្ទាត់ពុះ $\angle A$ នៃ $\triangle ABC$ កាត់ជ្រុង $BC$ ត្រង់ចំណុច $D$ ។ ពេលដែល $AB = 8$, $AC = 4$ និង $BD = 4$ រកប្រវែង $AD$ ។
  \item ពិចារណា $\triangle ABC$ មួយដែល $\angle A = 120^\circ$, $AB = 3$ និង $AC = 1$ ។ តាង $D$ ជាចំណុចប្រសព្វរវាង $BC$ និងកន្លះបន្ទាត់ពុះនៃ $\angle A$ រកប្រវែងនៃ $AD$ ។
  \item ឧបមាថាចតុកោណ $ABCD$ មួយ ចារឹកក្នុងរង្វង់មួយដែល $AB = 1$, $BC = 3$, $CD = 2$ និង $DA = 2$ ។ តាង $E$ ជាចំណុចប្រសព្វនៃអង្កត់ទ្រូង $AC$ និង $BD$ ។ \quad ក. រកផលធៀបនៃ $BE : ED$ \quad ខ. រកមុំ $BAD$ ។
  \item ឧបមាថា $0^\circ < a < 180^\circ < b < 360^\circ$ ។ កំណត់តម្លៃនៃចំនួនថេរ $a$ និង $b$ ដែល $\sin x + \sin(x + a) + \sin(x + b) = 0$ ចំពោះតម្លៃនៃ $x$ ។
  \item ឧបមាថា $0^\circ \le x \le 180^\circ$ និង $y = 3\sin 2x + a(\sin x + \cos x) + 1$ ដែល $a$ ជាចំនួនវិជ្ជមានថេរ ។
  \begin{enumerate}[a]
    \item យក $t = \sin x + \cos x$ រួចគណនា $f(x)$ ដោយប្រើ $t$
    \item រកដែនតម្លៃនៃ $t$
    \item រកតម្លៃអតិបរមានិងអប្បបរមានៃ $f(x)$ ។
  \end{enumerate}
  \item ពេលដែល $0^\circ \le x \le 180^\circ$ ដោះស្រាយសមីការខាងក្រោម ៖
  \begin{enumerate}[a]
    \item $\cos x + \cos(x + 30^\circ) + \cos(x + 60^\circ) = 0$
    \item $4\sqrt{2}\cos(x + 60^\circ)\cos(x + 45^\circ) = 3 + \sqrt{3}$ ។
  \end{enumerate}
  \item ឆ្លើយនឹងសំណួរខាងក្រោមអំពីត្រីកោណ $ABC$ ដែល $A = 60^\circ$ ។ \quad ក. រកដែននៃ $\sin B + \sin C$ \quad ខ. រកដែននៃ $\sin B\sin C$ ។
  \item រកតម្លៃនៃ $\sin\theta$ ចំពោះ $\theta = 18^\circ$ (ការណែនាំ ៖ $2\theta = 90^\circ - 3\theta$) ។
  \item ឧបមាថាចតុកោណ $ABCD$ មួយ ចារឹកក្នុងរង្វង់មួយដែល $AB = 2$, $BC = 6$, $CD = 4$ និង $B = 60^\circ$ ។
  \begin{enumerate}[a]
    \item រកប្រវែងនៃ $AC$ និង $AD$
    \item រកតម្លៃ $\sin C$
    \item ឧបមាថាបន្លាយនៃ $AB$ ជួបនិងបន្លាយ $CD$ ត្រង់ចំណុច $E$ ។ រកប្រវែងនៃ $AE$ និង $DE$ ។
  \end{enumerate}
\end{enumerate}
